\section{Examples}

\subsection*{Example 1: The Golden Ratio}

Let $\alpha = \varphi - 1 = (\sqrt{5}-1)/2 = [0; 1, 1, 1, \ldots]$.
The convergents are $p_k/q_k = F_{k-1}/F_k$ where $F_k$ is the $k$-th Fibonacci number:
\[
  \frac{0}{1},\; \frac{1}{1},\; \frac{1}{2},\; \frac{2}{3},\; \frac{3}{5},\;
  \frac{5}{8},\; \frac{8}{13},\; \ldots
\]
Since all partial quotients equal $1$, the gap lengths satisfy
$\ell_k = \|F_k \alpha\|$ and $L_k = \ell_{k-1}$, with
\[
  L_k / \ell_k \;\longrightarrow\; \varphi \approx 1.618\ldots \quad \text{as } k \to \infty.
\]
For example, with $n = 5$:

\begin{center}
\begin{tabular}{c|ccccc}
  $k$ & 1 & 2 & 3 & 4 & 5 \\ \hline
  $\{k\alpha\}$ & 0.618 & 0.236 & 0.854 & 0.472 & 0.090
\end{tabular}
\end{center}

Sorted: $0.090, 0.236, 0.472, 0.618, 0.854$.
Gaps: $0.146, 0.236, 0.146, 0.236, 0.236$. Wait---let us recompute cleanly.
With $\alpha = (\sqrt{5}-1)/2 \approx 0.6180$:
\[
  \{1\alpha\} \approx 0.618,\quad
  \{2\alpha\} \approx 0.236,\quad
  \{3\alpha\} \approx 0.854,\quad
  \{4\alpha\} \approx 0.472,\quad
  \{5\alpha\} \approx 0.090.
\]
Sorted order: $0.090 < 0.236 < 0.472 < 0.618 < 0.854$.
Gaps (including wrap-around):
\[
  0.146,\quad 0.236,\quad 0.146,\quad 0.236,\quad 0.236.
\]
Two distinct lengths: $\ell = 0.146 \approx \|5\alpha\| = |5\alpha - 3|$ and
$L = 0.236 \approx \|3\alpha\| = |3\alpha - 2|$.
Indeed $n = 5 = q_4$ is a Fibonacci number, so we expect exactly two lengths. $\checkmark$

\subsection*{Example 2: $\alpha = \sqrt{2} - 1$}

Here $\alpha = [0; 2, 2, 2, \ldots]$ and the convergents have Pell-number denominators:
$q_k = 1, 2, 5, 12, 29, 70, \ldots$ satisfying $q_{k+1} = 2q_k + q_{k-1}$.
The gap lengths at $n = q_k$ are
\[
  \ell_k = \|q_k \alpha\| = |q_k(\sqrt{2}-1) - p_k|,
\]
which decrease like $(\sqrt{2}-1)^k / q_k$.  The large partial quotient $a_k = 2$
means each long gap is split \emph{twice} before the next transition:
for $q_k < n \le q_k + q_{k-1}$ there are three gap lengths,
and for $q_k + q_{k-1} < n < q_{k+1}$ there are also three lengths but with
the second split already completed.

At $n = 7$ (between $q_2 = 5$ and $q_3 = 12$):
\[
  S_7(\alpha) \text{ has three gap lengths.}
\]
The lengths are approximately $0.172, 0.241, 0.413$, and indeed $0.172 + 0.241 = 0.413$. $\checkmark$

\subsection*{Example 3: A Rational $\alpha = 1/7$}

For $\alpha = 1/7$, the points $\{k/7\}$ for $k = 1,\ldots, 7$ are
$1/7, 2/7, 3/7, 4/7, 5/7, 6/7, 0$---equally spaced.
Exactly \textbf{one} gap length ($1/7$) for $n = 7$.

For $n = 5$ (not a multiple of $7$):
\[
  S_5(1/7) = \{1/7, 2/7, 3/7, 4/7, 5/7\}.
\]
Gaps: $1/7, 1/7, 1/7, 1/7, 3/7$.  Two distinct lengths ($1/7$ and $3/7$). $\checkmark$

The three-distance theorem applies to rational $\alpha$ too: with $\alpha = p/q$,
one always gets at most two gap lengths (never three), because the rational case
never produces a third length---the ``sum rule'' degeneracy doesn't arise.

\subsection*{Comparison Table}

\begin{center}
\renewcommand{\arraystretch}{1.3}
\begin{tabular}{l|c|c|c}
  $\alpha$ & CF & Denominators $q_k$ & Gap ratio $L/\ell$ \\
  \hline
  $(\sqrt{5}-1)/2$ & $[0;\overline{1}]$ & $1,2,3,5,8,13,\ldots$ & $\to \varphi$ \\
  $\sqrt{2}-1$ & $[0;\overline{2}]$ & $1,2,5,12,29,\ldots$ & $\to 1+\sqrt{2}$ \\
  $e - 2$ & $[0;1,2,1,1,4,1,1,6,\ldots]$ & $1,3,4,7,32,\ldots$ & varies \\
  $1/7$ & finite & $7$ only & $\le 2$ (rational)
\end{tabular}
\end{center}
